\paragraph{Strengths.}
The multi-stage pipeline design offers several advantages over end-to-end generation. The architecture itself was shaped by an early-stage informal experiment (see \textit{Architectural rationale} below): we initially tried prompting the model directly with raw commentary text and asking for a summary, and observed that the model consistently injected its own interpretive content alongside the commentators' positions. This failure mode---the inability to reliably distinguish \textit{what the commentators said} from \textit{what the model thinks they said}---drove the decision to introduce an explicit verbatim-quote extraction stage upstream of summarization, with token-set validation to enforce that the extracted quotes actually appear in the source. By extracting verbatim quotes before synthesis, the system creates an auditable intermediate representation that can be inspected and validated independently of the final BHT. The tiered commentator system ensures that the most authoritative voices dominate the output, while the constraint enforcement (word count, quote proportion, forbidden phrases) prevents many common failure modes of unconstrained LLM generation.

The iterative prompt refinement process (spanning multiple versions across three prompt families that shipped in the production run) proved essential. Early prompt versions produced BHTs that frequently exceeded word limits, included commentator names, or adopted list formatting. Later versions largely eliminated these issues, demonstrating the value of systematic prompt engineering with quantitative feedback loops.

\paragraph{Architectural rationale.}
The decision to extract verbatim quotes as a discrete pipeline stage, rather than relying on the model to summarize commentary directly, was made early in the project (mid-2023) on the basis of informal experimentation across roughly 10--20 key verses, evaluated and discussed at the project's weekly round-table review. Two failure modes were consistently observed when the model was given raw commentary and asked to produce a summary: (1) the model would supplement the commentators' actual positions with its own interpretive content, often stylistically indistinguishable from the source material; and (2) when asked to quote, the model would paraphrase rather than reproduce text verbatim, even when verbatim quotation was explicitly requested. This led to the two architectural commitments that define the production pipeline: a dedicated extraction stage that forces the model to produce text claimed to come from a specific commentator, and a token-set validation check that rejects any output containing more than two tokens not present in the source. Together, these allow the multi-stage pipeline to make a stronger claim than end-to-end summarization makes by default: every fact, opinion, and turn of phrase in the final BHT is intended to be traceable to a specific commentator's verbatim statement. We did not preserve formal records of these early experiments, and so do not present them as a quantitative ablation; we describe the reasoning here because it materially shaped the system's design.

\paragraph{Limitations.}
Several limitations warrant discussion. First, the entire BHT corpus was generated using GPT-3.5-turbo, which was among the most capable generally available models at the time of development; no regeneration has been performed on newer models because the pipeline was run once to produce the served corpus, and that one-time run's output quality remained acceptable while inference cost for newer models was substantially higher at the time. A full-scale regeneration and re-evaluation on more recent models remains future work. Second, the quality scoring framework, while multi-dimensional, relies heavily on embedding-based similarity, which may not capture all aspects of theological accuracy or interpretive nuance. Third, the quote-proportion constraint (30--90\% of BHT words drawn from source quotes, with a 70\% target), while ensuring grounding, can sometimes produce BHTs that read as stitched-together phrases rather than naturally flowing prose. Fourth, while we describe the early-stage experiments that motivated the multi-stage architecture (see \textit{Architectural rationale} above), we do not present a quantitative same-era zero-shot ablation against GPT-3.5-turbo, nor a frontier-model zero-shot baseline using a more recent system (e.g., GPT-4-class or Claude-class models). Such ablations would help isolate the marginal contribution of the choicest-extraction stage from the model itself, and remain useful future work; we note however that any zero-shot baseline would need to be specifically engineered to satisfy the source-traceability property that the production pipeline was designed to provide, so the comparison is best understood as one of summary quality conditional on giving up (or independently re-engineering) traceability---not of summary quality alone.

\paragraph{Source adherence vs.\ fluency.}
A fundamental tension exists between source adherence and fluency of the generated text. Higher quote proportions ensure adherence but can reduce coherence; lower proportions allow smoother prose but risk introducing LLM-generated content that may not reflect the commentators' actual positions. Our operating range of 30--90\% with a 70\% target represents a pragmatic balance, validated through iterative testing and user feedback; the observed corpus-level means (39.6\% for NT, 31.8\% for OT; see Table~\ref{tab:aggregate}) satisfy the 30\% floor and stay well below the 90\% ceiling, but consistently fall short of the 70\% target. In practice, produced BHTs paraphrase more than the target specifies; the gap between target and observed behavior is a known limitation and a natural candidate for follow-up tuning.

\paragraph{Prompt engineering insights.}
The evolution of prompts yielded several practical insights: (1)~explicit negative constraints (``do not mention commentators'') were more effective than positive-only instructions; (2)~specifying exact word count ranges reduced variance more than vague length instructions; (3)~the tiering language (``primarily,'' ``sparingly,'' ``minimally'') meaningfully influenced the tier distribution in generated outputs; and (4)~requesting numbered quotes in the extraction phase significantly improved downstream validation.

